\section{Einleitung}
\label{sec:Einleitung}

\subsection{Vorwort} 
\label{sub:Vorwort}
Im Jahre 2012 begann ich meine Ausbildung zum Fachinformatiker für Anwendungsentwicklung bei der \Abkuerzung{BTC}. Zuvor
habe ich eine schulische Ausbildung zum staatlich geprüfter TAI\footnote{Technischer Assistent für Informatik.} absolviert und
erlangte an derselben Schule die Fachhochschulreife. Seit Mitte Juli 2013 arbeite ich, fast durchgängig, in dem im
Kapitel~\ref{sub:Projektumfeld} beschriebenem Team.

Es liegt die Dokumentation des Projektes "`Erweiterung einer Webanwendung um Task Notifications"' vor. Dieses Projekt habe ich im
Rahmen der angestrebten IHK-Ausbildung bei der \Abkuerzung{BTC} durchgeführt. Es stellt einen möglichen Prozess der Erweiterung
einer bestehenden Webanwendung dar.

\subsection{Kurzporträt der Ausbildungsstätte} 
\label{sub:Kurzporträt der Ausbildungsstätte}
Die BTC Business Technology Consulting AG ist eines der führenden IT-Consulting-Unternehmen in Deutschland mit
Niederlassungen in der Schweiz, der Türkei, Polen und Japan. Das Unternehmen mit mehr als 1.600 Mitarbeitern und Hauptsitz in
Oldenburg ist Partner von SAP und Microsoft. \Abkuerzung{BTC} erzielte in 2013 einen Umsatz von 195,2 Mio. Euro.

\Abkuerzung{BTC} hat ein ganzheitliches, auf Branchen ausgerichtetes IT-Beratungsangebot und damit eine führende Position in den
Bereichen Energie, Telekommunikation, Öffentlicher Sektor, Industrie und Dienstleistungen. Das Angebot umfasst Beratung,
Systemintegration sowie Applikations- und Systemmanagement. Ein weiterer Fokus liegt auf energienahen Softwareprodukten.

Seit der Gründung im Jahre 2000 hat die \Abkuerzung{BTC} ein dichtes Netz an Geschäftsstellen aufgebaut, um IT-Consulting kundennah
erbringen zu können (Oldenburg, Berlin, Bremen, Hamburg, Leipzig, Mainz, Münster, Neckarsulm, Zürich, Poznań, Istanbul, Ankara,
Tokio). \citep{BTC}

\begingroup
\setlength\intextsep{0pt}

\subsection{Projektumfeld} 
\label{sub:Projektumfeld}
Das Projekt war im Entwicklerteam I des Bereiches "`Non-SAP Technologien"' der "`Business Unit Markets"' angesiedelt.
Dieses Team ist hauptsächlich für die Entwicklung des Kundenprojektes \Abkuerzung{TSO BSS} zuständig. Daraus resultiert unter
Anderem die Webanwendung \Abkuerzung{VMS}. Das Projekt ist als Teilprojekt der allgemeinen Weiterentwicklung des
\Abkuerzung{VMS} zu sehen. Im Folgenden gehe ich daher kurz auf die \Abkuerzung{TSO BSS} und das \Abkuerzung{VMS} ein.  

\begin{wrapfigure}[12]{r}[25pt]{0\textwidth}
\includegraphicsKeepAspectRatio{Bilder/Systemlandschaft.pdf}{0.6}
\caption{Systemlandschaft}
\label{app:Systemlandschaft} 
\end{wrapfigure}
Die \Abkuerzung{TSO BSS} gehört dem Sektor Energie an und ist ein Vertrags- und Kapazitätsmanagementsystem mit einer engen
Anbindung zu Vermarktungsplattformen. Die Webanwendung \Abkuerzung{VMS} dient der Verwaltung von Kundenverträgen. Zusätzlich
unterstützt das System sowohl in fachlichen als auch technischen Prozessen. Als Kommunikationszentrale spricht das
\Abkuerzung{VMS} mit nachgelagerten Systemen des \Abkuerzung{TSO} sowie mit externen Systemen, wie \zB Vermarktungsplattformen
(Siehe \Abbildung{app:Systemlandschaft}) \citep[vgl.][]{TSOBSS}.

\endgroup

\subsection{Motivation und Zielsetzung} 
\label{sub:Motivation und Zielsetzung}
\Abkuerzung{BTC} hat für das \Abkuerzung{VMS} eine Erweiterung identifiziert und dem Kunden vorgestellt. Der Kunde hat dann
die Umsetzung beauftragt.

Der Benutzer soll zu jeder Zeit auf einen Blick die Task-Anzahl, der ihm oder seiner Benutzerrolle zugewiesenen offenen Tasks,
sehen können. Hinsichtlich der Benutzerfreundlichkeit sollte diese Anzeige so in Form und Farbe gestaltet werden, dass
diese intuitiv und auch für Menschen mit Sehschwächen, wie Farbenblindheit, ohne Einschränkungen nutzbar ist \citep{Antrag}.

\subsection{Erklärung zum Dokument} 
\label{sub:Erklärung zum Dokument}
An dieser Stelle wird die Gestaltung des Dokuments erläutert. 
\begin{itemize}[noitemsep]
  \item Ein {[\textit{kursiver Text}]} weist auf eine Abkürzung hin, die im Abkürzungsverzeichnis zu finden ist.
  \item Ein {[\textbf{fettgedruckter Text}]}, mit Ausnahme von Überschriften, weist auf einen Fachbegriff hin, zu dem im Dokument
  eine Erklärung vorliegt. Ausgenommen sind Erklärungen in Fußnoten.
  \item Ein {[\texttt{Text in Schreibmaschinenschrift}]} weist auf eine Klasse oder Methode hin.
  \item Ein {[\textsf{Text in serifloser Schrift}]} weist auf sonstige Programmcode-Textbausteine hin.
  \item Liegt das Dokument in elektronischer Form vor, sind Abkürzungen und Quellenangaben zum jeweiligen Verzeichniseintrag
  verlinkt.
\end{itemize}

\subsection{Aufbau des Dokuments} 
\label{sub:Aufbau des Dokuments}
Einführend wird ein Überblick über den Aufbau der vorliegenden Projektdokumentation gegeben.

Im folgenden Kapitel wird auf die Rahmenbedingungen und Vorgehensweiße zur Durchführung des Projektes eingegangen.
Im dritten Kapitel wird die zugrundeliegende Webanwendung analysiert, Anforderungen an das Projekt definiert und die
Wirtschaftlichkeit dargelegt.
Die festgelegten Anforderungen werden im Kapitel Vier zu einem Konzept zusammengefasst.
Das anschließende fünfte Kapitel erklärt wie das Konzept umgesetzt wurde.
Nachfolgend wird im sechsten Kapitel aufgezeigt wie die Softwarequalität bewahrt wurde.
Abschließend wird die Projektdurchführung in Kapitel Sieben zusammengefasst betrachtet und ein persönliches Fazit gezogen.

